\chapter*{Synopsis}
% BEGIN CURATED SUBJECT INDEX
\index{operator!adjoint and self-adjoint}
\index{Hilbert space}
\index{topology}
\index{spectral theorem}
\index{spectrum!point}
\index{resolvent}
\index{von Neumann algebra@\vN\ algebra}
\index{Jost function}
\index{wave operator}
\index{partial-wave expansion}
% END CURATED SUBJECT INDEX
\addcontentsline{toc}{chapter}{Synopsis}

This book develops resonance theory as a calculational extension of ordinary
quantum mechanics.  Its starting point is not a non-Hermitian matrix but a
self-adjoint Hamiltonian, its scattering boundary data, and the global
connection problem of the associated differential equation.  Exact WKB
analysis supplies the constructive spine: formal Riccati solutions are Borel
resummed in sectors, Stokes jumps transport local bases, and the prohibited
incoming coefficient becomes an analytic function
$\Cincoming(E)$.  A zero of this function is simultaneously a
resolvent pole, a scattering pole, and---after an admissible change of
representation---a complex-scaled eigenvalue.  Its derivative gives the pole
residue and regularized normalization.

The book derives each layer needed to make that statement precise.  The
spectral theorem is developed through projection-valued measures and direct
integrals; domains and topologies are kept visible; \vN\ algebras and
their normal states are introduced as the natural language of spectral
sectors.  Two-body scattering is constructed from wave operators, resolvent
boundary values, partial waves, and Jost functions.  Zel'dovich
regularization, Berggren's Gaussian argument, complex scaling, the
Aguilar--Balslev--Combes theorem, rigged Hilbert spaces, and continuum level
density are then derived as different realizations or probes of one continued
connection datum.

The construction is enlarged in four directions.  Radial singularities add
Frobenius indices, Langer correction, and renormalized monodromy.  Weakly
bound nuclei add Jacobi coordinates, rearrangement channels, optical
potentials, continuum discretization, and CDCC; lithium projectiles and
beryllium-core breakup provide worked laboratories.  Multiple zeros produce
exceptional points, Jordan chains, self-orthogonality, and geometric
holonomy.  Black-hole horizons produce Regge--Wheeler-type connection
problems and quasinormal modes, for which exact WKB and complex scaling can be
compared in the same notation.

Long-range correlations mark a genuine boundary of the usual construction.
Two-body Coulomb scattering requires Coulomb-distorted rather than plane-wave
asymptotics.  Three charged particles require several overlapping asymptotic
sectors for which no single elementary global form is known.  Charged QED
states similarly require infrared dressing and may not belong to the vacuum
Fock representation.  These examples make topology, measure, representation,
and comparison dynamics into physical choices rather than mathematical
decoration.

Eight Stages introduce prerequisites only once.  Nine Laboratories then
reuse those canonical definitions to carry out full calculations.  The text
is intended to be self-contained at the level of derivation: references
record provenance, variants, and further results, but no essential step is
delegated to a cited paper.
